\documentclass[11pt,a4paper]{article}
\usepackage[utf8]{inputenc}
\usepackage[margin=1in]{geometry}
\usepackage{enumitem}
\usepackage{hyperref}
\usepackage{titlesec}

\titleformat{\section}{\Large\bfseries}{\thesection}{1em}{}
\titleformat{\subsection}{\large\bfseries}{\thesubsection}{1em}{}

\title{Project Presentation Guidelines}
\author{}
\date{}

\begin{document}

\maketitle

\section{General Guidelines}

In addition to the submitted report, each project must be presented by the authoring group, according to the following guidelines:

\begin{itemize}
    \item The presentation should be high level but sufficiently detailed information should be readily available to help answering questions from the audience.
    \item The duration of the presentation should be 10 minutes (groups of 1--2 students) or 15 minutes (groups of 3 students).
    \item At the end of the presentation there will be an extra 5--10 minutes of questions by anyone in the audience or two members of the course staff who are present. The questions from lecturer/TAs can be considered as an oral exam questions, and if answers to these questions reveal weak knowledge of the methods and workflow steps which should be part of the project, that can reduce the grade.
    \item Grading will be done by the two members of the course staff using standardized grading instructions.
\end{itemize}

\section{Specific Recommendations}

\begin{itemize}
    \item The first slide should include project's title and group members' names.
    \item The chosen statistical model(s), including observation model and priors, must be explained and justified.
    \item Make sure the font is big enough to be easily readable by the audience. This includes figure captions, legends and axis information.
    \item The last slide should be a summary or take-home-messages and include contact information or link to a further information. (The grade will be reduced by one if the last slide has only something like ``Thank you'' or ``Questions?'').
    \item In general, the best presentations are often given by teams that have frequently attended TA sessions and gotten feedback, so we strongly recommend attending these sessions.
\end{itemize}

\section{Details on Presentation Sessions in 2025}

\begin{itemize}
    \item Presentations will be given on campus.
    \item If you reserved a presentation slot but need to cancel, do it ASAP via the reservation link on MyCourses. If you cancel after the reservation link has closed, send a message to David Kohns (TA) on Zulip.
    \item As we have many presentation in each slot join the meeting in time. Late arrivals will lower the grade. Very late arrivals will fail the presentation and can present in period III.
    \item It is easiest if just one from the group shares the slides, but it is expected that all group members present some part of the presentation orally.
    \item Presentation time is 10 min for 1--2 person groups and 15 min for 3 person groups.
    \item Time limit is strict. It's good idea to practice the talk so that you get the timing right. Staff will announce 2 min and 1 min left and time ended. Going overtime reduces the grade.
    \item After the presentation there will be 5 min for questions, answers, and feedback.
    \item Each student has to come up with at least one question during the session. Students can ask more questions.
    \item Staff will ask further questions (kind of oral exam).
\end{itemize}

\section{Grading Criteria}

Grading of the project presentation takes into account following things in the order of importance. You can get feedback for all parts.

\subsection{Students Show Understanding of Essential Concepts}
\begin{itemize}[label=--]
    \item TAs will ask clarifying questions to check these
\end{itemize}

\subsection{Inclusion of Required Workflow Steps}

\subsection{Clarity of Problem, Model and Prior Description}

\subsection{Focus of the Presentation}
\begin{itemize}[label=--]
    \item Focus on the most relevant points in the limited time
\end{itemize}

\subsection{Overall Clarity and Coherence}
\begin{itemize}[label=--]
    \item How well the slides and presentation flows
    \item Effective presentation of the results
\end{itemize}

\subsection{Slide Design and Visual Aspects}
\begin{itemize}[label=--]
    \item E.g., use of figures, font size, number of digits
    \item Effective use of figures
\end{itemize}

\subsection{Lack of Typos, Grammatical and Spelling Errors}

\subsection{Oral Presentation}
\begin{itemize}[label=--]
    \item Clarity of oral presentation, rhythm and timing, rate of speech
    \item Clarity of voice projection and appropriate volume
    \item Presenting style created interest and was engaging
    \item Presence (e.g., speaking to the audience/camera)
    \item Possible distracting mannerisms
    \item \textbf{Note:} This has minor effect, and involuntary aspects (e.g., ticks) do not affect grading
\end{itemize}

\subsection{Timing}
\begin{itemize}[label=--]
    \item Point reduction if overtime
    \item Fair amount of time for all group members
\end{itemize}

\subsection{Bonus: Going Beyond the Expectations}
\begin{itemize}[label=--]
    \item You may get a bonus point for going beyond the explicit requirements
\end{itemize}

\section{Self-Evaluation}

Students will also self-evaluate their project. Instructions for this will be provided.

\end{document}